% !TEX program = pdflatex
\documentclass[11pt,a4paper]{article}

% =========================================================================
% GUÍA DE TRABAJO DEL INTEGRANTE
% Proyecto: CasaScore (AlquilaPuno)
% Gestión de Proyectos TICS (EST315) -- Semestre 2026-II -- Team Sobrevivientes
% =========================================================================

% --- PAQUETES DE CONFIGURACIÓN ---
\usepackage[utf8]{inputenc}
\usepackage[T1]{fontenc}
\usepackage[spanish,es-tabla]{babel}
\usepackage{geometry}
\geometry{
    top=3.6cm,
    bottom=2.3cm,
    left=2.2cm,
    right=2.2cm,
    % PARTE MODIFICABLE: headheight debe ser mayor que la altura del logo más alto (2.56cm = 73pt).
    headheight=80pt,
    headsep=15pt,
    footskip=20pt
}

\usepackage{graphicx}
\usepackage{fancyhdr}
\usepackage[table]{xcolor}
\usepackage{titlesec}
\usepackage{booktabs}
\usepackage{tabularx}
\usepackage{array}
\usepackage{amsmath}
\usepackage{enumitem}
\usepackage{microtype}
\usepackage{amssymb}
% --- RECURSOS VISUALES: GRÁFICOS Y CÓDIGO DE COLORES ---
\usepackage{tikz}
\usetikzlibrary{positioning,calc}
\definecolor{GrisFila}{RGB}{246, 248, 251}
\definecolor{CuadCerca}{RGB}{150, 30, 30}
\definecolor{CuadSatisf}{RGB}{180, 130, 20}
\definecolor{CuadInfo}{RGB}{10, 45, 90}
\definecolor{CuadMonit}{RGB}{110, 115, 125}
\newcommand{\etiqueta}[2]{{\setlength{\fboxsep}{2pt}\colorbox{#1}{\textcolor{white}{\scriptsize\bfseries #2}}}}
\newcommand{\cerca}{\etiqueta{CuadCerca}{Gestionar de cerca}}
\newcommand{\satisf}{\etiqueta{CuadSatisf}{Mantener satisfecho}}
\newcommand{\informado}{\etiqueta{CuadInfo}{Mantener informado}}
\newcommand{\monit}{\etiqueta{CuadMonit}{Monitorear}}
\usepackage{needspace}
\usepackage{longtable}
\definecolor{Rama1}{RGB}{10, 45, 90}
\definecolor{Rama2}{RGB}{40, 95, 150}
\definecolor{Rama3}{RGB}{180, 130, 20}
\definecolor{Rama4}{RGB}{30, 110, 60}
\definecolor{Rama5}{RGB}{110, 60, 140}
\definecolor{RaciR}{RGB}{150, 30, 30}
\definecolor{RaciA}{RGB}{10, 45, 90}
\definecolor{RaciC}{RGB}{180, 130, 20}
\definecolor{RaciI}{RGB}{140, 145, 155}
\definecolor{NivAlto}{RGB}{30, 110, 60}
\definecolor{NivMedio}{RGB}{180, 130, 20}
\definecolor{NivBajo}{RGB}{150, 30, 30}
\usetikzlibrary{patterns}
\definecolor{SevBaja}{RGB}{30, 110, 60}
\definecolor{SevMedia}{RGB}{180, 130, 20}
\definecolor{SevAlta}{RGB}{150, 30, 30}
\usetikzlibrary{arrows.meta,shapes.geometric,fit,backgrounds}
\usepackage{listings}
\lstset{basicstyle=\ttfamily\scriptsize, frame=single, rulecolor=\color{AzulUNA!40}, backgroundcolor=\color{GrisFila}, columns=fullflexible, keepspaces=true, breaklines=true, xleftmargin=4pt, xrightmargin=4pt, literate={á}{{\'a}}1 {é}{{\'e}}1 {í}{{\'i}}1 {ó}{{\'o}}1 {ú}{{\'u}}1 {ñ}{{\~n}}1 {Á}{{\'A}}1}

\usepackage{hyperref}

% Logos: funciona a 2 o 3 niveles de profundidad dentro de docs/latex/
\graphicspath{{../assets/img/}{../../assets/img/}{../../../assets/img/}}

% --- PALETA INSTITUCIONAL FIEI / UNA PUNO ---
\definecolor{AzulUNA}{RGB}{10, 45, 90}
\definecolor{DoradoFIEI}{RGB}{180, 130, 20}
\definecolor{GrisOscuro}{RGB}{50, 50, 50}
\definecolor{AzulClaro}{RGB}{228, 236, 246}

% --- CONFIGURACIÓN DE LISTAS ---
\setlist{topsep=2pt, itemsep=2pt, parsep=1pt}

% --- COLUMNAS DE TABLA ---
\newcolumntype{L}{>{\raggedright\arraybackslash}X}
\newcolumntype{P}[1]{>{\raggedright\arraybackslash}p{#1}}
\newcolumntype{C}[1]{>{\centering\arraybackslash}p{#1}}
\newcolumntype{Q}[1]{>{\raggedleft\arraybackslash}p{#1}}
\newcolumntype{R}{>{\raggedleft\arraybackslash}X}
\renewcommand{\arraystretch}{1.25}

% --- FORMATO DE ENLACES ---
\hypersetup{
    colorlinks=true,
    linkcolor=AzulUNA,
    urlcolor=AzulUNA,
    citecolor=AzulUNA,
    pdftitle={Guía De Trabajo Del Integrante -- CasaScore (AlquilaPuno)},
    pdfauthor={Team Sobrevivientes}
}

% --- FORMATO DE TÍTULOS ---
\titleformat{\section}{\color{AzulUNA}\normalfont\large\bfseries}{\thesection.}{0.7em}{}[\color{AzulUNA}\titlerule]
\titleformat{\subsection}{\color{AzulUNA}\normalfont\normalsize\bfseries}{\thesubsection.}{0.5em}{}
\titlespacing*{\section}{0pt}{9pt}{4pt}
\titlespacing*{\subsection}{0pt}{6pt}{2pt}

% --- ENCABEZADO Y PIE DE PÁGINA (fancyhdr) ---
\pagestyle{fancy}
\fancyhf{}

% ===== TAMAÑO DE LOS LOGOS (PARTE MODIFICABLE) =====
% logo_fiei.png  : height=1.45cm (lado izquierdo)
% logo_equipo.png: height=2.56cm (lado derecho; 25 % más grande porque la imagen tiene más borde blanco)
% Para cambiar el tamaño, modifique solo el valor height=... de cada \includegraphics.
% Si el logo más alto supera 2.6cm, aumente también headheight y top en la configuración de página.
\fancyhead[L]{%
    \includegraphics[height=1.45cm,keepaspectratio]{logo_fiei.png}%
}
\fancyhead[C]{%
    \scriptsize \textbf{\textcolor{AzulUNA}{UNIVERSIDAD NACIONAL DEL ALTIPLANO DE PUNO}}\\[2pt]
    \tiny \textcolor{GrisOscuro}{FACULTAD DE INGENIERÍA ESTADÍSTICA E INFORMÁTICA}\\[2pt]
    \scriptsize \textbf{\textcolor{AzulUNA}{Gestión de Proyectos TICS (EST315) \textbar\ Team Sobrevivientes}}\\[2pt]
    \footnotesize \textbf{\textcolor{DoradoFIEI}{GUÍA DE TRABAJO DEL INTEGRANTE}}%
}
\fancyhead[R]{%
    \includegraphics[height=2.56cm,keepaspectratio]{logo_equipo.png}%
}

\renewcommand{\headrulewidth}{0pt}
\renewcommand{\headrule}{}

\fancyfoot[L]{\scriptsize \textcolor{GrisOscuro}{CasaScore (AlquilaPuno)}}
\fancyfoot[C]{\small \thepage}
\fancyfoot[R]{\scriptsize \textcolor{GrisOscuro}{Puno, Perú -- Semestre 2026-II}}
\renewcommand{\footrulewidth}{0.4pt}

\begin{document}
\shorthandoff{<>}
\begin{center}
    {\Large \textbf{\textcolor{AzulUNA}{GUÍA DE TRABAJO DEL INTEGRANTE}}}\\[3pt]
    {\normalsize \textbf{Ruth Karina Apaza Solis -- Analista de Planificación}}\\[4pt]
    \footnotesize \textbf{Fecha:} jueves 05/11/2026 (semana 11) \quad \textbar\quad \textbf{Repositorio:} \nolinkurl{xClaIDx/ProyectoTics} \quad \textbar\quad \textbf{Encarga:} Clyde Neil Paricahua Pari
\end{center}

\vspace{0.1cm}

\section{Propósito}
El equipo trabaja desde ahora en el repositorio \texttt{ProyectoTics}. Esta guía indica cómo subir a ese repositorio los documentos que elaboraste, uno por uno y con su Pull Request, y cómo llevar tus tarjetas en el tablero \textit{CasaScore -- Tablero híbrido}.

\vspace{0.1cm}
\noindent\begin{tabularx}{\textwidth}{@{}P{3.6cm} L@{}}
\toprule
\textbf{Repositorio:} & \nolinkurl{https://github.com/xClaIDx/ProyectoTics} \\
\textbf{Tablero:} & \nolinkurl{https://github.com/users/xClaIDx/projects/1} \\
\textbf{Tu usuario:} & \texttt{R-Karina-A-Solis} \\
\textbf{Revisor:} & Clyde (\texttt{xClaIDx}) revisa y aprueba cada Pull Request \\
\textbf{Plazo:} & PR 1 a 7 hasta el viernes 06/11; PR 8 el mismo viernes, después del corte de valor ganado \\
\bottomrule
\end{tabularx}

\section{Lo que Recibes}
\begin{enumerate}[leftmargin=*]
    \item \textbf{El archivo \texttt{paquete\_ruth.zip}}, que Clyde te envía por WhatsApp. Al descomprimirlo verás la carpeta \nolinkurl{paquete_ruth/docs/latex/} con tus 14 documentos (\texttt{.tex} y \texttt{.pdf}) ordenados en las mismas carpetas que tienen en el repositorio.
    \item \textbf{La invitación al repositorio} \texttt{ProyectoTics}, que llega a tu correo de GitHub.
    \item \textbf{Tus tarjetas en el tablero}, ya asignadas a tu usuario.
\end{enumerate}

\Needspace{12\baselineskip}
\section{Antes de Empezar (una sola vez)}
\begin{enumerate}[leftmargin=*]
    \item Acepta la invitación desde el correo (\textbf{View invitation} $\rightarrow$ \textbf{Accept invitation}) o en \nolinkurl{github.com/xClaIDx/ProyectoTics/invitations}.
    \item Abre el repositorio y comprueba que ya existen las carpetas \nolinkurl{docs/latex/casascore/}, \nolinkurl{docs/latex/horus_sentinel/} y \nolinkurl{docs/latex/00_comun/}. Las sube Clyde antes que tú; si no las ves, avísale.
    \item Descomprime \texttt{paquete\_ruth.zip} en tu escritorio.
    \item Abre el tablero, filtra con \texttt{assignee:R-Karina-A-Solis} y revisa tus tarjetas.
\end{enumerate}

\Needspace{12\baselineskip}
\section{Cómo se Simula el Historial}
Cada documento conserva la fecha que tiene escrita; no se modifica ningún archivo. GitHub registrará el día real en que lo subes. Para que se sepa cuándo debió entregarse, el mensaje del commit y el título del Pull Request terminan con la fecha prevista, por ejemplo:

\begin{lstlisting}
docs(casascore-viabilidad): viabilidad y business case [fecha prevista 10/09/2026]
\end{lstlisting}

\Needspace{24\baselineskip}
\section{Tus Ocho Pull Requests}
Súbelos en este orden. En cada uno, las carpetas se toman de \nolinkurl{paquete_ruth/docs/latex/} y se suben a la misma ruta dentro de \nolinkurl{docs/latex/} del repositorio.

\vspace{0.1cm}
{\small
\rowcolors{2}{white}{GrisFila}
\noindent\begin{tabularx}{\textwidth}{@{}C{0.6cm} P{4.05cm} L C{1.8cm}@{}}
\toprule
\rowcolor{AzulClaro}
\textbf{PR} & \textbf{Rama} & \textbf{Carpetas que subes} & \textbf{Fecha prevista} \\
\midrule
1 & \nolinkurl{docs/casascore-viabilidad} & \nolinkurl{casascore/03_viabilidad_business_case} y \nolinkurl{horus_sentinel/03_viabilidad_business_case} & 10/09/2026 \\
2 & \nolinkurl{docs/casascore-edt} & \nolinkurl{casascore/06_edt} y \nolinkurl{horus_sentinel/06_edt} & 18/09/2026 \\
3 & \nolinkurl{docs/casascore-plan-equipo} & \nolinkurl{casascore/07_plan_equipo} y \nolinkurl{horus_sentinel/07_plan_equipo} & 18/09/2026 \\
4 & \nolinkurl{docs/casascore-cronograma} & \nolinkurl{casascore/08_cronograma} y \nolinkurl{horus_sentinel/08_cronograma} & 18/09/2026 \\
5 & \nolinkurl{docs/casascore-presupuesto} & \nolinkurl{casascore/09_presupuesto} y \nolinkurl{horus_sentinel/09_presupuesto} & 18/09/2026 \\
6 & \nolinkurl{docs/horus-control} & \nolinkurl{horus_sentinel/16_control_proyecto} y \nolinkurl{00_comun/C7_comparacion_financiera} & 21/09/2026 \\
7 & \nolinkurl{docs/casascore-control} & \nolinkurl{casascore/16_control_proyecto} & 02/10/2026 \\
8 & \nolinkurl{docs/casascore-desempeno} & \nolinkurl{casascore/21_informe_desempeno} & 06/11/2026 \\
\bottomrule
\end{tabularx}
}

\vspace{0.1cm}
\noindent\textbf{Mensajes de commit:}

\begin{lstlisting}
PR 1  docs(casascore-viabilidad): viabilidad y business case de ambos proyectos [fecha prevista 10/09/2026]
PR 2  docs(casascore-edt): EDT y diccionario de ambos proyectos [fecha prevista 18/09/2026]
PR 3  docs(casascore-equipo): plan del equipo de ambos proyectos [fecha prevista 18/09/2026]
PR 4  docs(casascore-cronograma): cronograma de ambos proyectos [fecha prevista 18/09/2026]
PR 5  docs(casascore-presupuesto): presupuesto de ambos proyectos [fecha prevista 18/09/2026]
PR 6  docs(horus-control): control al corte de la semana 4 y comparacion financiera [fecha prevista 21/09/2026]
PR 7  docs(casascore-control): control del proyecto al corte de la semana 6 [fecha prevista 02/10/2026]
PR 8  docs(casascore-desempeno): informe de desempeno al corte de la semana 11 [fecha prevista 06/11/2026]
\end{lstlisting}

\Needspace{30\baselineskip}
\section{Paso a Paso desde el Navegador}
Es la forma más sencilla: no necesitas instalar nada. El ejemplo es el PR 1; los demás se hacen igual cambiando la rama, las carpetas y el mensaje.

\subsection{Crear la rama}
\begin{enumerate}[leftmargin=*]
    \item Entra a \nolinkurl{github.com/xClaIDx/ProyectoTics}.
    \item Haz clic en el botón de la rama, que dice \textbf{main}, arriba a la izquierda de la lista de archivos.
    \item Escribe \nolinkurl{docs/casascore-viabilidad} y elige \textbf{Create branch docs/casascore-viabilidad from main}. El botón ahora muestra el nombre de tu rama.
\end{enumerate}

\subsection{Subir las carpetas}
\begin{enumerate}[leftmargin=*]
    \item Con tu rama seleccionada, entra a \texttt{docs} $\rightarrow$ \texttt{latex} $\rightarrow$ \texttt{casascore}.
    \item Pulsa \textbf{Add file} $\rightarrow$ \textbf{Upload files}.
    \item Desde tu escritorio, arrastra la carpeta completa \nolinkurl{paquete_ruth/docs/latex/casascore/03_viabilidad_business_case} al recuadro. GitHub conserva la carpeta con sus dos archivos.
    \item Abajo, en \textbf{Commit changes}, escribe el mensaje del PR 1. Deja marcada la opción \textbf{Commit directly to the docs/casascore-viabilidad branch} y pulsa \textbf{Commit changes}.
    \item Vuelve a \nolinkurl{docs/latex/}, entra a \texttt{horus\_sentinel} (sin cambiar de rama) y repite con la carpeta \nolinkurl{horus_sentinel/03_viabilidad_business_case}. En el mensaje agrega al final ``(Horus Sentinel)''.
\end{enumerate}

\subsection{Abrir el Pull Request}
\begin{enumerate}[leftmargin=*]
    \item Aparecerá un aviso amarillo con el botón \textbf{Compare \& pull request}; púlsalo. Si no aparece, entra a la pestaña \textbf{Pull requests} $\rightarrow$ \textbf{New pull request} y elige tu rama.
    \item Verifica arriba: \textbf{base: main} $\leftarrow$ \textbf{compare: docs/casascore-viabilidad}.
    \item Título: el mismo del mensaje de commit. En la descripción copia esta plantilla:
\end{enumerate}

\begin{lstlisting}
Fecha prevista: 10/09/2026 (semana 3)
Que incluye: Viabilidad y Business Case de CasaScore y de Horus Sentinel (.tex y .pdf)
Tarjeta del tablero: "Subir mis documentos a ProyectoTics (8 PR)" - PR 1
Horas: 0.5
\end{lstlisting}

\begin{enumerate}[leftmargin=*, resume]
    \item A la derecha, en \textbf{Reviewers}, elige \texttt{xClaIDx}.
    \item Pulsa \textbf{Create pull request}. No lo fusiones tú: Clyde lo revisa, lo aprueba y pulsa \textbf{Merge pull request}.
\end{enumerate}

\Needspace{14\baselineskip}
\subsection{Después de cada Pull Request}
\begin{enumerate}[leftmargin=*]
    \item En el tablero abre la tarjeta \textit{Subir mis documentos a ProyectoTics (8 PR)}, marca la casilla del PR y suma las horas en \textit{Horas reales}.
    \item Si Clyde deja observaciones, entra al Pull Request, corrige en la misma rama (\textbf{Add file} $\rightarrow$ \textbf{Upload files} reemplaza el archivo) y responde el comentario.
    \item Cuando Clyde fusione el Pull Request, la rama se puede borrar con el botón \textbf{Delete branch}.
    \item Para el siguiente Pull Request, crea la nueva rama siempre desde \textbf{main}.
\end{enumerate}

\Needspace{20\baselineskip}
\section{Tus Tarjetas en el Tablero}
{\small
\rowcolors{2}{white}{GrisFila}
\noindent\begin{tabularx}{\textwidth}{@{}P{4.4cm} C{1.9cm} L@{}}
\toprule
\rowcolor{AzulClaro}
\textbf{Tarjeta} & \textbf{Estado} & \textbf{Qué haces} \\
\midrule
Subir mis documentos a ProyectoTics (8 PR) & Por hacer & Muévela a \textit{En Progreso} al empezar. Marca cada PR al abrirlo y pásala a \textit{Hecho} cuando Clyde fusione el PR 8. \\
Montar el tablero híbrido & En Progreso & En la vista \textit{Tablero}, abre el menú de la columna \textit{En Progreso} $\rightarrow$ \textbf{Set limit} y pon 4; en \textit{En Revision} pon 3. Luego pásala a \textit{Hecho}. \\
Corte de valor ganado semana 11 & En Progreso & El viernes 06/11 calcula PV, EV, AC, SPI y CPI con la vista \textit{Control semanal}; ese corte es el que sube el PR 8. \\
Seguridad por fila (RLS) & En Progreso & Apoya a Eliseo con la seguridad por fila de Supabase hasta el 18/11. \\
Viabilidad y Acta de Constitución; Cronograma y presupuesto de Horus; Control de Horus y comparación financiera; Planificación base & Hecho & Ya están cerradas. Cuando fusionen tu PR, pega el enlace del Pull Request como comentario en la tarjeta que le corresponde. \\
\bottomrule
\end{tabularx}
}

\vspace{0.1cm}
\noindent\textbf{Horas:} anota en \textit{Horas reales} el tiempo de cada sesión el mismo día, en decimales (media hora $=$ 0.5).

\Needspace{20\baselineskip}
\section{Alternativa desde la Terminal}
Si prefieres Git, haz lo mismo con estos comandos. El ejemplo es el PR 1; la carpeta \texttt{paquete\_ruth} debe estar en tu escritorio.

\begin{lstlisting}
git clone https://github.com/xClaIDx/ProyectoTics.git
cd ProyectoTics
git checkout main && git pull
git checkout -b docs/casascore-viabilidad
cp -r ~/Desktop/paquete_ruth/docs/latex/casascore/03_viabilidad_business_case docs/latex/casascore/
cp -r ~/Desktop/paquete_ruth/docs/latex/horus_sentinel/03_viabilidad_business_case docs/latex/horus_sentinel/
git add docs/latex/casascore/03_viabilidad_business_case docs/latex/horus_sentinel/03_viabilidad_business_case
git commit -m "docs(casascore-viabilidad): viabilidad y business case de ambos proyectos [fecha prevista 10/09/2026]"
git push -u origin docs/casascore-viabilidad
\end{lstlisting}

\noindent Después abre el Pull Request en GitHub como se indica en la sección anterior.

\Needspace{16\baselineskip}
\section{Si Algo Sale Mal}
{\small
\rowcolors{2}{white}{GrisFila}
\noindent\begin{tabularx}{\textwidth}{@{}P{5.0cm} L@{}}
\toprule
\rowcolor{AzulClaro}
\textbf{Situación} & \textbf{Solución} \\
\midrule
No aparece \textbf{Add file} o dice que no tienes permiso & Aún no aceptaste la invitación al repositorio; acéptala y recarga la página. \\
Subiste los archivos a \texttt{main} sin querer & No hagas nada más y avisa a Clyde; él lo corrige. \\
Subiste a una carpeta equivocada & En la misma rama abre el archivo, menú \textbf{\ldots} $\rightarrow$ \textbf{Delete file}, y vuelve a subirlo en la carpeta correcta. \\
El Pull Request dice \textit{This branch has conflicts} & Avisa a Clyde por WhatsApp con el enlace del Pull Request. \\
No encuentras tu tarjeta en el tablero & Quita los filtros y escribe \texttt{assignee:R-Karina-A-Solis} en la barra de filtro. \\
\bottomrule
\end{tabularx}
}

\end{document}
